%%=============================================================================
%% Frontendarchitectuur en gebruikersinterface
%%=============================================================================

\chapter{Frontendarchitectuur en gebruikersinterface}%
\label{ch:frontendarchitectuur}

Dit hoofdstuk bespreekt de frontend. De applicatie is gebouwd met Expo React Native. De frontend is opgesplitst in screens, hooks, components, API-modules, utilities, validation en styles. Die structuur maakt het makkelijker om gedrag te hergebruiken en schermen overzichtelijk te houden.

\section{Schermen, hooks en componenten}

Screens bepalen welke onderdelen op het scherm verschijnen. Hooks beheren gedrag zoals data laden, zoeken, filteren en pagineren. Components zijn herbruikbare UI-blokken. API-modules bevatten de communicatie met de backend.

\begin{figure}[H]
  \centering
  \includegraphics[width=\linewidth]{Frontend_architecture.png}
  \caption[Frontendarchitectuur]{Schema van de frontendstructuur met navigation, screens, hooks, herbruikbare componenten en ondersteunende modules.}
  \label{fig:frontend-architecture}
\end{figure}

\clearpage

\section{Dashboardlogica}

De dashboards bewaren de ruwe tickets als bron van waarheid. Daarna worden zoekresultaten, filters, statusgroepen, tellingen en pagina's afgeleid. Dit vermindert de kans op inconsistente lijsten. Als de ruwe ticketlijst verandert, worden de afgeleide resultaten opnieuw berekend.

Voor de adminpagina ziet de pipeline er conceptueel zo uit:

\begin{quote}
klanten en tickets $\rightarrow$ zoeken $\rightarrow$ filters $\rightarrow$ statusgroepen $\rightarrow$ paginatie
\end{quote}

Voor de klantpagina is de pipeline eenvoudiger:

\begin{quote}
tickets $\rightarrow$ zoeken en filteren $\rightarrow$ statusgroepen $\rightarrow$ paginatie
\end{quote}

\begin{figure}[H]
  \centering
  \includegraphics[width=0.95\linewidth,height=0.268\textheight,keepaspectratio]{Filter_tickets.png}
  \caption[Ticketfilters]{Ticketoverzicht met zoekveld en filters voor status, categorie en onderwerp.}
  \label{fig:ticket-filters}
\end{figure}

\section{Paginatie en hergebruik}

Paginatie is gecentraliseerd in \texttt{usePagination}. De hook is generiek gemaakt, zodat hij niet alleen voor tickets werkt, maar ook voor gesprekken en klantenlijsten. Daardoor hoeft dezelfde paginalogica niet op meerdere plaatsen opnieuw geschreven te worden.

\begin{listing}[H]
\begin{minted}[fontsize=\scriptsize,baselinestretch=0.9]{typescript}
export function usePagination<T>(
  items: T[],
  { pageSize, initialPage = "first", resetKey }: UsePaginationOptions,
) {
  const totalPages = Math.max(1, Math.ceil(items.length / pageSize));
  const [currentPage, setCurrentPage] = useState(1);

  const pageItems = useMemo(
    () => items.slice(startIndex, endIndex),
    [endIndex, items, startIndex],
  );

  return { pageItems, currentPage, totalPages };
}
\end{minted}
\caption[Generieke paginatiehook]{Vereenvoudigd voorbeeld van de generieke paginatiehook.}
\label{lst:usepagination}
\end{listing}

\section{Ticketgroepen}

De klant- en adminpagina gebruiken statusgroepen voor \texttt{New}, \texttt{Open} en \texttt{Closed}. \texttt{TicketGroupSection} toont een groep tickets en gebruikt intern dezelfde paginatiehook. Elke groep heeft eigen paginatie. Bladeren in gesloten tickets verandert dus niet de pagina van nieuwe tickets.

\section{Conversatieweergave}

In het ticketdetail wordt de conversatie weergegeven met berichten aan verschillende kanten. Adminberichten en klantberichten krijgen ook een andere achtergrondkleur. Dit lijkt meer op chatapplicaties zoals WhatsApp of Messenger en maakt een lange conversatie sneller scanbaar.

\begin{figure}[H]
  \centering
  \includegraphics[width=0.95\linewidth,height=0.64\textheight,keepaspectratio]{Ticket_conversation.png}
  \caption[Ticketconversatie]{Ticketdetail met chatbubbels, visueel onderscheid tussen klant en admin en afbeeldingspreview.}
  \label{fig:ticket-conversation}
\end{figure}
